% Isolated statement and pre-existing inputs only; no receiving proof, computed residue basis or calculation output.
% Actor preparing isolation: GPT-6 Astra (OpenAI), ChatGPT Pro mode, inherited reasoning effort.
\section*{Isolated Paper V upper-extension target}

The target statement is copied below. Independently establish that the displayed
connection represents the supplied nonsplit complex flat extension, including its
invariant line, and determine the stated contraction scalar. In particular derive
the equivalence of flat bundles that proves existence of a common monodromy
intertwiner; interval residual containment is not an input.

\begin{theorem}\label{night:P1.15}
The rank-six flat extension, as a complex flat bundle with its invariant
line, has the following explicit algebraic connection model.  Set
$\omega_k=x^kdx/y$, $0\leq k\leq3$, and use the quotient frame
\[
(w_2,w_1,w_0,w_{-1},w_{-2})
 =(\omega_{01},\omega_{02},\omega_{03},\omega_{13},\omega_{23}),
\qquad \omega_{12}=-3\omega_{03}.
\]
These are meromorphic frames of
$(L^2,L,\mathcal O,L^{-1},L^{-2})$, where $L=T^2$ on the root
orbifold of orders $5,4$.  If a frame is written as a row of vectors,
its derivative is right multiplication by
\[
\Gamma_6=\begin{pmatrix}0&b\\0&\Gamma_W\end{pmatrix},\qquad
b=\frac1{\lambda-\lambda_*}
 \left(-\frac{25}{6\lambda^2},-\frac{125}{48\lambda},0,1,0\right),
\tag{1.15.1}
\]
where
\[
\resizebox{\textwidth}{!}{$\displaystyle
\Gamma_W=\frac1{3125\lambda+256}
\begin{pmatrix}
-8(625\lambda+48)/\lambda&-8&10\lambda&75\lambda^2/2&0\\
-125/(2\lambda)&-4(625\lambda+48)/\lambda&8&0&75\lambda^2/2\\
300/\lambda^2&-375/\lambda&0&-48&60\lambda\\
120/\lambda^3&0&125/(2\lambda)&4(625\lambda+48)/\lambda&-8\\
0&120/\lambda^3&50/\lambda^2&-125/(2\lambda)&8(625\lambda+48)/\lambda
\end{pmatrix}.$}
\tag{1.15.2}
\]
This model specifies a canonical algebraic splitting, denoted
$\sigma_b$.  The compatible contraction used for the real parameter is
$\alpha=(625\pi i/128)\alpha_{\rm alg}$, where
$\alpha_{\rm alg}(w_0)=4/3$ and it vanishes on $w_1,w_2$.
\end{theorem}

\section*{Pre-existing inputs and boundary conventions}

The base is the root curve of orders four at $\lambda=0$ and five at
$\lambda=\infty$, with one omitted nodal point
$a=-256/3125$. Set
$L=\mathcal O((4/5)[\infty]-(3/4)[0])$ and
$W=L^2\oplus L\oplus\mathcal O\oplus L^{-1}\oplus L^{-2}$,
in exactly the five-column order of the target.
Invariant Laurent-series descent gives
$\pi_*L^j=\mathcal O_{\mathbb P^1}(\lfloor4j/5\rfloor+
\lfloor-3j/4\rfloor)$.

The quotient connection is the displayed $\Gamma_W$; its identification with
the integral quotient local system is supplied by P1.2, P1.7 and the independent
polynomial Gauss--Manin reduction. The supplied rank-six flat local system
has invariant line $\mathbb C e$ and this quotient. Its canonical extension on
the root curve is a holomorphic extension of $W$ by $\mathcal O e$.
It is regular on both finite root discs and has only a logarithmic pole at $a$,
where its nilpotent residue has square zero and rank three.
Here $\lambda=u^4$ at zero, and the regular quotient frame is
$(u^6w_2,u^3w_1,w_0,u^{-3}w_{-1},u^{-6}w_{-2})$.
At infinity $\lambda=v^{-5}$ and the regular quotient frame is
$(v^{-8}w_2,v^{-4}w_1,w_0,v^4w_{-1},v^8w_{-2})$.
The invariant section $e$ itself is regular and nowhere zero.
The target uses row frames, $\nabla\mathcal F=\mathcal F\Gamma$.

The supplied cohomology monodromies, for fixed based positive meridians,
are the matrices
\[
A=\begin{pmatrix}
1&0&0&1&-1&1\\0&1&-1&-1&1&0\\0&1&0&-1&0&0\\
0&2&-1&-1&0&1\\0&2&0&-1&0&0\\0&1&1&-1&0&0
\end{pmatrix},\quad
B=\begin{pmatrix}
1&0&0&0&0&0\\0&0&-1&0&1&0\\0&0&-1&0&0&1\\
0&0&-1&0&0&0\\0&1&-1&0&0&0\\0&0&-1&1&0&0
\end{pmatrix}.
\]
P1.1 gives the invariant line $\mathbb C e_0$ and the absence of a
nonzero common invariant vector for the dual representation. These assertions
may also be checked directly from the matrices. The quotient is identified in
the same fixed meridians with $\Gamma_W$ as above.

The rank-three lifted Hodge plane projects to
$P=\langle w_2,w_1,w_0\rangle\subset W$.
The algebraic contraction $\alpha_{\rm alg}:P\to\mathcal O$ sends $w_0$ to
$4/3$ and $w_1,w_2$ to zero. The actual marked cusp contraction is determined
on the one-dimensional kernel of the quotient residue restricted to $P$ by
$N_V(\widetilde v)=\alpha(v)e$, where logarithmic monodromy is $2\pi i$
times the nilpotent residue. The given marking of the invariant line is $e_0$.

For a fundamental solution with identity initial value at $\lambda=1$, prove
existence of a constant invertible $X$ that intertwines both actual monodromies
with $A,B$ and satisfies $Xe_0=e_0$. Check its uniqueness by the simultaneous
centralizer, or identify any obstruction. The only target is this exact flat
extension and its normalization; no numerical period calculation is required.
